\section*{Appendix: Hyperparameter Search Methodology}
\label{app:hyperparam}

\subsection*{Algorithm and Environment}

We tune vanilla SAC (no expert guidance) on the \texttt{Plane} environment with
target and initial altitudes drawn uniformly from $[3{,}000, 8{,}000]$\,m and
episodes capped at $10{,}000$ steps.
The optimisation objective is the mean episodic reward evaluated over 25 test
episodes, logged every $10{,}000$ environment steps.

\subsection*{Search Algorithm}

We use \textbf{Optuna's Tree-structured Parzen Estimator (TPE)}~\cite{bergstra2011tpe},
persisted in a SQLite database (\texttt{hp\_results/optuna.db}) for full
resume-safety.
After an initial random phase of roughly 10 trials, TPE fits two kernel density
estimates $p(\mathbf{x} \mid \text{good})$ and $p(\mathbf{x} \mid \text{bad})$
from completed trials and proposes the next configuration by maximising the ratio
$p(\mathbf{x} \mid \text{good}) / p(\mathbf{x} \mid \text{bad})$.

\subsection*{Two-Phase Protocol}

Because the SAC training loop is a monolithic JAX-JIT call without mid-training
callbacks, standard asynchronous successive halving (ASHA) cannot prune running
trials.
We instead adopt a two-phase design that recovers most of ASHA's sample
efficiency:

\begin{description}
  \item[Phase~1 — Exploration.]
    Each of the $N = 40$ TPE trials is trained for
    $3 \times 10^{5}$ timesteps over \textbf{5 independent seeds}.
    The trial score is the final (last-logged) value of the mean episodic reward,
    averaged across seeds.

  \item[Phase~2 — Confirmation.]
    The top-$K$ configurations from Phase~1 are re-trained for
    $10^{6}$ timesteps over \textbf{15 independent seeds}.
    The trial score is the mean episodic reward averaged over all evaluation
    points in the last $2 \times 10^{5}$ timesteps and across all seeds,
    matching the evaluation window used in the final experiments.
\end{description}

\subsection*{Search Space}

Table~\ref{tab:hp_space} lists the full search space.
Continuous ranges marked \emph{log} are sampled uniformly on a log scale.

\begin{table}[h]
\centering
\caption{Hyperparameter search space for SAC on \texttt{Plane}.}
\label{tab:hp_space}
\begin{tabular}{llll}
\toprule
\textbf{Hyperparameter} & \textbf{Type} & \textbf{Range / Values} & \textbf{Scale} \\
\midrule
Actor learning rate       & Continuous & $[10^{-4},\, 10^{-3}]$    & log   \\
Critic learning rate      & Continuous & $[10^{-4},\, 10^{-3}]$    & log   \\
Entropy coefficient lr    & Continuous & $[10^{-4},\, 10^{-3}]$    & log   \\
Discount factor $\gamma$  & Continuous & $[0.97,\, 0.999]$          & linear\\
Polyak coefficient $\tau$ & Continuous & $[5\times10^{-4},\, 5\times10^{-2}]$ & log \\
Initial entropy coeff $\alpha_0$ & Continuous & $[0.05,\, 2.0]$      & log   \\
Target entropy per dim    & Continuous & $[-4.0,\, -0.2]$           & linear\\
Batch size                & Categorical & $\{256,\, 512\}$          & —     \\
Network width             & Categorical & $\{256,\, 512\}$          & —     \\
Gradient clipping         & Categorical & $\{\text{none},\, 0.5,\, 1.0\}$ & — \\
\bottomrule
\end{tabular}
\end{table}

All actor and critic networks are two-hidden-layer MLPs with ReLU activations and
equal width $w \in \{256, 512\}$, i.e.\ architecture $(w, \texttt{relu}, w,
\texttt{relu})$.
Gradient clipping, when enabled, clips by global $\ell_2$ norm.

\subsection*{Compute Budget}

Phase~1 consumes at most $40 \times 5 \times 3\times10^{5} = 6\times10^{7}$
environment steps.
Phase~2 (top-$K=5$) adds at most $5 \times 15 \times 10^{6} = 7.5\times10^{7}$
steps.
Trials are parallelised across available GPUs; each trial subprocess receives a
dedicated GPU via \texttt{CUDA\_VISIBLE\_DEVICES}.

% ============================================================
\section*{Appendix: Hyperparameter Tables}
\label{app:hptables}

All experiments share the same environment: \texttt{Plane} with target and
initial altitudes in $[3{,}000, 8{,}000]$\,m, episodes of at most $10{,}000$
steps, and $1 \times 10^{6}$ total training timesteps.
Reward is evaluated over 25 test episodes after every $10{,}000$ environment
steps.

% ------------------------------------------------------------
\subsection*{SAC Hyperparameters}

Table~\ref{tab:sac_tuned} lists the hyperparameters selected by the Bayesian
search described above; Table~\ref{tab:sac_fixed} lists the remaining
hyperparameters that were held fixed across all experiments.

\begin{table}[h]
\centering
\caption{Tuned SAC hyperparameters (Phase~2 best configuration).}
\label{tab:sac_tuned}
\begin{tabular}{lll}
\toprule
\textbf{Hyperparameter} & \textbf{Value} & \textbf{Notes} \\
\midrule
Actor learning rate       & $3.41 \times 10^{-4}$ & log-uniform search \\
Critic learning rate      & $2.55 \times 10^{-4}$ & log-uniform search \\
Entropy coeff.\ lr        & $6.13 \times 10^{-4}$ & log-uniform search \\
Discount $\gamma$         & $0.9805$               & linear search       \\
Polyak $\tau$             & $8.36 \times 10^{-3}$ & log-uniform search  \\
Initial ent.\ coeff.\ $\alpha_0$ & $1.207$         & log-uniform search  \\
Target entropy per dim    & $-0.268$               & linear search       \\
Batch size                & $256$                  & categorical         \\
Network width             & $256$                  & categorical         \\
Gradient clip ($\ell_2$)  & $1.0$                  & categorical         \\
\bottomrule
\end{tabular}
\end{table}

\begin{table}[h]
\centering
\caption{Fixed SAC hyperparameters (not tuned, shared by all algorithms).}
\label{tab:sac_fixed}
\begin{tabular}{lll}
\toprule
\textbf{Hyperparameter} & \textbf{Value} & \textbf{Notes} \\
\midrule
Replay buffer size        & $1{,}000{,}000$        & transitions \\
Learning starts           & $10{,}000$             & steps before first update \\
Number of environments    & $1$                    & \\
Reward scale              & $1.0$                  & \\
Number of critics         & $2$                    & double Q-learning \\
Network architecture      & $(w, \texttt{ReLU}, w, \texttt{ReLU})$ & $w=256$ \\
Evaluation episodes       & $25$                   & per checkpoint \\
Log / eval frequency      & $10{,}000$ steps       & \\
Training timesteps        & $1 \times 10^{6}$      & \\
Independent seeds         & $100$                  & \\
\bottomrule
\end{tabular}
\end{table}

% ------------------------------------------------------------
\subsection*{EDGE Hyperparameters}

Table~\ref{tab:ege_hp} reports the hyperparameters for \textbf{EDGE}
(\texttt{ege\_decay\_050}), the main result.
EDGE gates between the expert action and the learned policy action using
$p_{\text{expert}} = \sigma\!\left(\Delta Q / (\tau\,|Q|)\right)$, where
$\Delta Q = Q(s, \pi^*) - Q(s, \pi_\theta)$ and $|Q|$ is a running scale
estimate.
The gating probability decays linearly from $\varepsilon_0$ to $0$ over the
first $f_{\mathrm{decay}}$ fraction of training.

\begin{table}[h]
\centering
\caption{EDGE hyperparameters (main result: \texttt{ege\_decay\_050}).}
\label{tab:ege_hp}
\begin{tabular}{lll}
\toprule
\textbf{Hyperparameter} & \textbf{Value} & \textbf{Notes} \\
\midrule
Gating mode               & Probabilistic (sigmoid) & $p_{\text{expert}} = \sigma(\Delta Q / \tau |Q|)$ \\
Initial expert probability $\varepsilon_0$ & $0.5$ & \\
Decay fraction $f_{\mathrm{decay}}$ & $0.50$ & fraction of total training \\
Gating temperature $\tau$ & $1.0$  & \\
Expert buffer pre-population & $0$  & buffer fills organically \\
Expert replay mix fraction & $0.0$ & \\
Online BC auxiliary loss  & disabled & \\
Critic updates per step   & $1$    & \\
Policy update start       & $2{,}000$ steps & \\
Alpha update start        & $2{,}000$ steps & \\
\bottomrule
\end{tabular}
\end{table}

% ------------------------------------------------------------
\subsection*{IBRL and Residual Policy Hyperparameters}

Table~\ref{tab:baselines_hp} lists the hyperparameters specific to the two
additional baselines.
IBRL~\cite{hu2023ibrl} gates deterministically: the expert action is taken
whenever $Q(s, \pi^*) > Q(s, \pi_\theta)$ (argmax, no decay).
Residual RL~\cite{johannink2019residual} outputs a residual $a_\delta$ that is
added to the (frozen) expert action before environment execution:
$a = \mathrm{clip}(a_{\pi^*} + a_\delta,\, -1, 1)$.

\begin{table}[h]
\centering
\caption{Hyperparameters for IBRL and Residual RL baselines.}
\label{tab:baselines_hp}
\begin{tabular}{lll}
\toprule
\textbf{Hyperparameter} & \textbf{IBRL} & \textbf{Residual RL} \\
\midrule
Gating / action mode         & Argmax (deterministic) & $\mathrm{clip}(a_{\pi^*} + a_\delta,\,-1,1)$ \\
Expert gating decay          & none ($f_{\mathrm{decay}}=0$) & — \\
Expert buffer pre-population & $0$            & $0$ \\
Expert replay mix fraction   & $0.0$          & $0.0$ \\
Online BC auxiliary loss     & disabled       & disabled \\
Critic updates per step      & $1$            & $1$ \\
Policy update start          & $2{,}000$ steps & $10{,}000$ steps \\
Alpha update start           & $2{,}000$ steps & $10{,}000$ steps \\
\bottomrule
\end{tabular}
\end{table}

% ============================================================
\section*{Appendix: Statistical Tests}
\label{app:stats}

\subsection*{Overview}

All pairwise comparisons are made between EDGE (\texttt{ege\_decay\_050}, the
main result) and each of the three baselines: vanilla SAC, Residual RL, and
IBRL.
Using a single reference method eliminates logical constraints between the
tested hypotheses, so we apply the Holm–Bonferroni step-down procedure for
multiplicity control.
Each family therefore contains $K = 3$ comparisons and controls the
family-wise error rate at $\alpha = 0.05$.

\subsection*{Test Statistic: Inter-Quartile Mean (IQM)}

The Inter-Quartile Mean (IQM) trims the bottom and top $25\%$ of the seed
distribution and averages the remainder.
For a sorted sample $x_{(1)} \le \cdots \le x_{(n)}$:
\begin{equation}
  \mathrm{IQM}(x) =
    \frac{1}{h - l}
    \sum_{i=l}^{h-1} x_{(i)},
  \quad
  l = \lfloor 0.25\,n \rfloor,\;
  h = \lceil  0.75\,n \rceil.
  \label{eq:iqm}
\end{equation}
We operate on per-seed scalar summaries (described below) so that each seed
contributes exactly one observation to the IQM.
With $100$ seeds the IQM discards the $25$ lowest- and $25$ highest-performing
seeds, averaging the central $50$.

\subsection*{Test Families}

Two separate families of tests are conducted on the
\textbf{Normalised Expert Advantage} metric, defined as
\begin{equation}
  \mathrm{NEA}_t =
    \frac{\bar{r}_t^{\pi} - \bar{r}_t^{*}}{\mathrm{MAX\_SCORE} - \bar{r}_t^{*}},
\end{equation}
where $\bar{r}_t^{\pi}$ is the policy's mean episodic reward at step $t$,
$\bar{r}_t^{*}$ is the expert's mean episodic reward, and
$\mathrm{MAX\_SCORE}$ is the theoretical maximum.
Higher values indicate a policy closer to expert-level performance; $\mathrm{NEA}=0$
means the policy matches the expert.

\paragraph{Asymptote.}
For each seed, the IQM of all evaluation checkpoints in the \emph{last}
$200{,}000$ environment steps is computed, yielding one scalar per seed.
This measures long-run performance.

\paragraph{Steps to threshold.}
For each seed, the first step at which NEA permanently crosses a
threshold $y^*$ is recorded.
Seeds that never cross $y^*$ within the $10^6$-step budget are
\emph{budget-censored}: they receive the value $10^6$ (RMST-style imputation).
Thresholds $y^* \in \{0.0,\, 0.2,\, 0.4\}$ measure how quickly each method
closes $0\%$, $20\%$, and $40\%$ of the gap to the expert.

\subsection*{Permutation Test}

For a comparison between groups $A$ (EDGE) and $B$ (baseline) with per-seed
observations $\mathbf{x}_A$ and $\mathbf{x}_B$, the observed test statistic is
\begin{equation}
  T_{\mathrm{obs}} = \mathrm{IQM}(\mathbf{x}_A) - \mathrm{IQM}(\mathbf{x}_B).
\end{equation}
Under the null hypothesis $H_0$: the two groups are exchangeable, we pool
$\mathbf{x} = [\mathbf{x}_A; \mathbf{x}_B]$ and draw $B = 10{,}000$
permutations.
For each permutation $\pi^{(b)}$ the statistic
$T^{(b)} = \mathrm{IQM}(\pi^{(b)}_{1:|A|}) - \mathrm{IQM}(\pi^{(b)}_{|A|+1:})$
is computed.
The two-sided Monte Carlo p-value is
\begin{equation}
  p_{\mathrm{raw}} =
    \frac{1}{B}\sum_{b=1}^{B}
    \mathbf{1}\!\left[|T^{(b)}| \ge |T_{\mathrm{obs}}|\right].
\end{equation}
The random number generator is seeded at $42$ for reproducibility.

\subsection*{Holm–Bonferroni Correction}

Let $p_{(1)} \le p_{(2)} \le p_{(3)}$ be the sorted raw p-values for the three
comparisons and let $H_{(k)}$ denote the corresponding hypothesis.
The Holm–Bonferroni adjusted p-values are
\begin{equation}
  \tilde{p}_{(k)} = \max_{j \le k}\,\min\!\left(1,\; p_{(j)} \cdot (K - j + 1)\right),
  \quad k = 1, 2, 3.
\end{equation}
$H_{(k)}$ is rejected if $\tilde{p}_{(k)} < 0.05$.
This step-down procedure controls the FWER at $\alpha = 0.05$ under arbitrary
dependence between tests.

\subsection*{Results}

Tables~\ref{tab:stats_asymptote} and \ref{tab:stats_thresholds} report the
outcomes.
$\Delta\,\mathrm{IQM} = \mathrm{IQM}_{\mathrm{EDGE}} - \mathrm{IQM}_{\mathrm{baseline}}$;
for the asymptote a positive value means EDGE achieves higher normalized expert
alignment; for steps-to-threshold a negative value means EDGE reaches the target
level faster.
Only Holm–Bonferroni adjusted p-values are reported.

\begin{table}[h]
\centering
\caption{IQM permutation test + Holm–Bonferroni: asymptote
         (IQM of Normalised Expert Advantage over last $200{,}000$ steps,
         $100$ seeds, $10{,}000$ permutations, $\alpha=0.05$).
         Higher NEA is better; $\Delta\,\mathrm{IQM} > 0$ means EDGE achieves
         greater expert alignment at convergence.}
\label{tab:stats_asymptote}
\begin{tabular}{lrrr}
\toprule
\textbf{EDGE vs} & \textbf{$\Delta\,\mathrm{IQM}$} & \textbf{Adj.\ $p$} & \textbf{Reject?} \\
\midrule
SAC baseline  & $+0.12$ & $<0.0001$ & Yes \\
Residual RL   & $+0.04$ & $0.0002$  & Yes \\
IBRL          & $+0.03$ & $0.0091$  & Yes \\
\bottomrule
\end{tabular}
\end{table}

\begin{table}[h]
\centering
\caption{IQM permutation test + Holm–Bonferroni: steps to threshold
         (Normalised Expert Advantage, budget-censored at $10^6$ steps,
         $100$ seeds, $10{,}000$ permutations, $\alpha=0.05$).
         Negative $\Delta\,\mathrm{IQM}$ means EDGE reaches the NEA threshold
         earlier (fewer steps).
         ``Censored'' shows budget-censored seeds as EDGE\,/\,baseline out of $100$.}
\label{tab:stats_thresholds}
\begin{tabular}{llrrll}
\toprule
\textbf{Threshold $y^*$} & \textbf{EDGE vs} & \textbf{$\Delta\,\mathrm{IQM}$ (steps)} & \textbf{Adj.\ $p$} & \textbf{Reject?} & \textbf{Censored} \\
\midrule
\multirow{3}{*}{$0.0$}
  & SAC baseline & $-161{,}300$ & $0.1068$  & No  & $6\,/\,3$  \\
  & Residual RL  & $+27{,}000$  & $0.7508$  & No  & $6\,/\,1$  \\
  & IBRL         & $-148{,}400$ & $0.2018$  & No  & $6\,/\,8$  \\
\midrule
\multirow{3}{*}{$0.2$}
  & SAC baseline & $-211{,}300$ & $0.0039$  & Yes & $9\,/\,6$  \\
  & Residual RL  & $-94{,}300$  & $0.2313$  & No  & $9\,/\,2$  \\
  & IBRL         & $-139{,}600$ & $0.1878$  & No  & $9\,/\,10$ \\
\midrule
\multirow{3}{*}{$0.4$}
  & SAC baseline & $-177{,}600$ & $<0.0001$ & Yes & $10\,/\,6$  \\
  & Residual RL  & $-107{,}600$ & $0.0498$  & Yes & $10\,/\,3$  \\
  & IBRL         & $-69{,}000$  & $0.2823$  & No  & $10\,/\,10$ \\
\bottomrule
\end{tabular}
\end{table}
